\section{Open questions and the boundary of the classification}\label{sec:open}

Two questions remain open in the enumeration, and one limitation of
the classification deserves a precise statement.

\paragraph{Rigidity.}
The first question is the continuous-convolution rigidity assertion
\[
 0\ne A\in\CC[\rho],\qquad
 A\mid\int_0^\rho A(t)A(\rho-t)\dd t
 \quad\Longrightarrow\quad A=c\rho^{\deg A}.
\]
Through the polynomial Laplace correspondence used in the proof of
Theorem~\ref{thm:elliptic-count}, it says that, after fixing the pole
at zero, any rational one-pole solution of \eqref{eq:main} vanishing at
infinity would have to be \(c_\ast z^{-p}\), with operator
\(P=D^p\).
Algebraic and \(\ell\)-adic arguments, together with exact
computations, prove it in the range
\(\deg A\le572\), corresponding to \(2\le p\le573\), and in infinite
arithmetic families of higher degrees.
Rigidity in arbitrary degree remains a conjecture.
By Proposition~\ref{prop:general}, it would establish both
enumeration formulas, as counts with multiplicity, at every order.

\paragraph{Simplicity.}
The second question is whether every point of \(\Xp\), the
rational-exponential matching system, is simple. The elliptic
simplicity assertion is generic in the lattice and is derived from
this one in Theorem~\ref{thm:elliptic-count}.
At a solution \(S_A=AC_A\), the tangent map of the system is
\begin{equation}\label{eq:tangent}
 J_A(h)=\remA\bigl(2(A\star h)-C_Ah\bigr),\qquad \deg h<d.
\end{equation}
The question is whether this map always has zero kernel.
The prime-index argument proves this after reduction modulo
\(\ell\), where the map becomes multiplication by a polynomial
coprime to \(\overline A\); a uniform replacement is missing at
composite indices.
One might expect simplicity to follow from the uniqueness of the
Laurent expansion in Theorem~\ref{thm:complete}. It does not: Laurent
uniqueness concerns a fixed operator, whereas \eqref{eq:tangent} allows
the operator to vary with the profile.

\begin{conjecture}\label{conj:count}
Continuous-convolution rigidity holds in every degree, and the
matching Jacobian \eqref{eq:tangent} is invertible at every point
of \(\Xp\). Consequently there are exactly \(\binom{2p-1}{p-1}\)
rational-exponential pairs, all distinct, for every \(p\ge2\).
\end{conjecture}

The conjecture would also imply \(M_p\) simple elliptic pairs
on a generic fixed lattice at every order, \(F_p\) distinct pairs
in the two pure subfamilies, and that the rational-exponential waves of
every purely dissipative equation are fronts
(Theorem~\ref{thm:symmetric}).
Simplicity on every lattice is false already at order two;
\eqref{eq:ks-collisions} describes the full exceptional set
at order three.

\paragraph{Mixed equations with \(p\) even.}
The hypothesis of Theorem~\ref{thm:complete} is a separate issue. For
even \(p\) it requires a purely dispersive equation, and it cannot
simply be dropped. The known example comes from the Fisher equation.
Its travelling waves of nonzero speed satisfy \eqref{eq:main} with
\(p=2\) and a mixed operator, which is also the profile equation of
KdV--Burgers. At the
special speed of \citet{AblowitzZeppetella1979} it has a one-parameter
family of solutions, given in \citet{Kudryashov2010},
\citet[Lemma~4.4]{ConteNgWu2015} and
\citet{McCueElHachemSimpson2021}; see also
\citet{KudryashovFisher2005}. In our normalization it is
\[
 v(z)=-12\lambda^2e^{2\lambda z}
       \wp(e^{\lambda z}-1;0,g_3),\qquad
 P(D)=D^2-5\lambda D+6\lambda^2,\qquad\lambda\ne0.
\]
Substitution follows from
\((D^2-5\lambda D+6\lambda^2)[t^2U(t)]
=\lambda^2t^4U''(t)\), \(t=e^{\lambda z}\).
For \(g_3=0\) and \(\lambda=1\) this is the front in the first row of
\eqref{eq:kdvb-pairs}. For generic \(g_3\ne0\) it is globally
meromorphic and lies outside \(\Wclass\).
Further examples of meromorphic solutions outside \(\Wclass\)
are discussed by \citet[Remark~1.11]{NgWu2019}.
Thus a classification for mixed operators with \(p\) even must allow a
larger class of functions than \(\Wclass\).

For mixed operators of even degree \(p\ge4\) we do not know whether
meromorphic solutions outside \(\Wclass\) occur.
Methods for constructing elliptic and rational-exponential
solutions without the finiteness property are developed by
\citet{Demina2022}.
